\section{Representations}
\label{sec:ch07-representations}
\index{representations|(}%

A representation is how the router names one instance of an entity to a subgraph
that has never seen it. It is a JSON object with a \code{\_\_typename} and the
fields of a \code{@key}, and nothing else:

\begin{minted}{json}
{ "__typename": "Product", "id": "1" }
\end{minted}

That is the whole idea. Reviews holds no products, has no product table and
cannot list one. What it can do is take that object and answer questions about
the thing it names.

\index{federation!entities field@\code{\_entities} field}%
The field that accepts them is \code{\_entities}, and you can send it by hand.
This is the request the router makes for the chapter's query, typed into Postman
instead:

\begin{minted}{graphql}
query ($representations: [_Any!]!) {
  _entities(representations: $representations) {
    ... on Product {
      __typename
      reviews { rating body }
    }
  }
}
\end{minted}

\begin{minted}{json}
{
  "representations": [
    { "__typename": "Product", "id": "1" },
    { "__typename": "Product", "id": "2" },
    { "__typename": "Product", "id": "3" }
  ]
}
\end{minted}

Three things about that document are worth naming before the answer.

\index{representations!in variables}%
The representations travel in a \emph{variable}, not inlined into the query
text. That is what lets the router compile one document and reuse it for a batch
of any size, and it is why the query the router sends for three products is
character for character the query it would send for three hundred.

The selection set is inside an inline fragment on \code{Product}, because
\code{\_entities} returns \code{[\_Entity]}, and \code{\_Entity} is a union. Every
selection on it has to say which member it is for. A query touching two entity
types in one call carries two fragments.

And \code{\_Any} is a scalar, which means the representations are not validated
against an input type. The specification puts the rules in prose instead: a
representation must carry \code{\_\_typename}, and it must carry every field named
in the \code{@key} it is standing in for~\autocite{apollo2026subgraphspec}. Get
that wrong and you find out at runtime.

\subsection*{The answer is positional}

\begin{minted}[fontsize=\footnotesize]{json}
{"data":{"_entities":[
  {"__typename":"Product","reviews":[
     {"rating":5,"body":"Survived a house move."},
     {"rating":3,"body":"Arrived with a chipped leg."}]},
  {"__typename":"Product","reviews":[
     {"rating":4,"body":"Firmer than it looks."}]},
  {"__typename":"Product","reviews":[]}
]}}
\end{minted}

Three representations in, three entities out, in the same order. There is no
identifier in the response tying an answer back to its representation, and there
does not need to be: answer $n$ belongs to representation $n$. That is the
contract, and it is also the reason a subgraph must never reorder or deduplicate
the list it was handed.

Product 3 has no reviews and its slot holds an empty list. It was still asked
about, and it still cost a slot in the batch.

\index{HotChocolate!EntitiesResolver@\code{EntitiesResolver}}%
Inside HotChocolate, \code{EntitiesResolver.ResolveAsync} is what runs this. Read
at tag \code{16.6.0}, commit \code{8fea46e}: it clones the resolver context once
per representation, starts every reference resolver as its own task, and awaits
them all afterwards. The representations in one call are resolved concurrently,
not in a loop. If one of them throws, the error is reported at path
\code{\_entities[n]} and that slot becomes null while the rest of the list
survives, which is the behaviour the specification's nullable
\code{[\_Entity]} return type exists to permit.

\subsection*{The subgraph will answer for things that do not exist}

There is no product 9. Ask anyway:

\begin{minted}{json}
{ "representations": [ { "__typename": "Product", "id": "9" } ] }
\end{minted}

\begin{minted}{json}
{ "data": { "_entities": [ { "id": "9", "reviews": [] } ] } }
\end{minted}

No error and no null: a perfectly ordinary \code{Product} with no reviews on it.

This is chapter~\ref{ch:the-federation-model}'s ownership rule arriving on the
wire, and it is sharper here than it was in the abstract. Reviews' reference
resolver is three lines long:

\begin{minted}{csharp}
[ReferenceResolver]
public static Product ResolveById(string id)
    => new() { Id = id };
\end{minted}

It takes a key and builds an object from it. It does not check that the product
exists, because it has no way to check: Catalog is the authority on which
products exist and Reviews has never been introduced to it. Asking would mean a
subgraph calling another subgraph behind the router's back, which is the one
thing federation is arranged to avoid.

In this direction it is harmless, because the router only ever sends keys that
Catalog just handed it. It stops being harmless the moment a client can reach
\code{\_entities} directly, and
chapter~\ref{ch:security-hardening} is where the subgraph stops being
reachable from outside the cluster.

\subsection*{A representation the subgraph cannot resolve}

\code{Review} has no \code{@key}, so it is not in the \code{\_Entity} union and
there is no reference resolver for it anywhere. Send one:

\begin{minted}{json}
{ "representations": [ { "__typename": "Review", "id": "r1" } ] }
\end{minted}

\begin{minted}{json}
{
  "errors": [ { "message": "Unexpected Execution Error", "path": ["_entities"] } ],
  "data": null
}
\end{minted}

HTTP 200, and a message that tells you nothing at all. I am printing it because
it is what you will get, not because it is any good.

\code{EntityResolver\_NoResolverFound} is what HotChocolate threw inside, and
that exception knows exactly what went wrong. What survives to the client is the
masked form, and the mask is the right default for a production subgraph even
when it is infuriating in development.

\index{federation!diagnosing a missing key}%
So learn the shape rather than the message. \enquote{Unexpected Execution Error}
at path \code{\_entities}, on a subgraph that is otherwise healthy, means the
router asked for a type this subgraph cannot resolve by key. Three causes, in
the order I check them: the type has no \code{@key} here; it has one but no
reference resolver; or it has both and is unreachable from any root field, which
is section~\ref{sec:ch07-what-a-subgraph-owes}'s trap wearing a different hat.
\code{\_service} answers all three in one request.

\subsection*{Verify it in Postman}

\index{Postman!federated-wire collection}%
Every response in this section and the last one is a request in the companion
repository's federated-wire collection, and the collection is worth running
rather than reading. Both files live in \code{postman/}:

\begin{minted}{text}
federated-wire.postman_collection.json
federated-wire.local.postman_environment.json
\end{minted}

The environment carries the three URLs. Import both and run the folder.

Nine requests, thirty-seven assertions. Three of them go straight to a subgraph
on its own port, before any router is involved: what Catalog publishes, what
Reviews publishes, and the \code{\_entities} call above with its three
representations. Two more are the failure cases: the product that does not exist
and the type that cannot be resolved. The last three go through the router.

\begin{minted}{text}
npx newman run postman/federated-wire.postman_collection.json \
    -e postman/federated-wire.local.postman_environment.json
\end{minted}

\index{Postman!assertions on a plan}%
The assertions worth stealing are on the plan rather than the data. One request
checks that the second fetch is a \code{BatchEntity} sent to \code{reviews},
that its \code{dependsOnFetchIds} names the first fetch, and that the first
fetch's query contains \code{\_\_typename} and \code{id}. Another sends a query
that crosses no boundary and asserts the opposite: one fetch, and neither field
present. A test on the answer would pass in both cases and tell you nothing.

\index{representations|)}%
